\chapter{幾何分布}
    \section{無記憶性からの導出}
    \begin{proof}
        \quad\par
        確率変数$X$を、ベルヌーイ試行を繰り返して初めて成功するまでに要する試行回数とする。
        無記憶であるとはすなわち「試行開始から$n$回失敗したという事象の下、それからさらに$k\in\naturalNumbers$回失敗する条件付き確率」が「試行開始から連続で$k$回失敗する確率」に等しいということである。
        特に$k=1$としたときの条件から次の式が成り立つ。
        \begin{align*}
            \frac{P(X>n+1)}{P(X>n)} &= P(X>1) \\
            P(X>n+1) &= P(X>1)P(X>n) \\
            P(X>n) - P(X>n+1) &= P(X>n) - P(X>1)P(X>n) \\
            P(X=n+1) &= (1-P(X>1))P(X>n) = P(X=1)(1-P(X\leq n))
        \end{align*}
        これより
        \begin{align*}
            P(X=2) &= P(X=1)(1 - P(X=1)) = p(1-p) \quad (p \coloneqq P(X=1)) \\
            P(X=3) &= P(X=1)(1 - P(X=1) - P(X=2)) = p[1-p-p(1-p)] = p(1-p)^2 \\
            P(X=4) &= P(X=1)(1 - P(X=1) - P(X=2) - P(X=3)) \\
            &= p[1-p-p(1-p)-p(1-p)^2] = p(1-p)^3 \\
            &\vdots \\
            P(X=n) &= p(1-p)^{n-1}
        \end{align*}
    \end{proof}